\chapter{Drawing the Screen}

The last two chapters were about the server: a process that holds a session
with the plant and pushes values to whoever is connected. This one is about
what's connected: a page in a browser that receives those values and draws the
plant on the screen.

The browser is a strange place to run a program. The code isn't installed. It
travels over the network every time someone opens the page, so its size is paid
for by every visitor. It runs on a machine you don't control, so anything it
contains can be read by anyone who opens the page. And it can only talk to the
network under rules the browser enforces, whatever the server thinks. Most of
what makes front-end work different from server work comes from those three
facts, and from one older problem: keeping what's on the screen in agreement
with what the program knows.

\section{A page you change by hand}
\label{sec:by-hand}

\begin{history}{1995--2006 \textperiodcentered{} From a few lines to jQuery}
In May 1995, Brendan Eich wrote the first version of JavaScript in ten days at
Netscape, and it shipped in Navigator 2.0, in betas from the end of that year. It was made for small
things: checking a form before sending it, swapping an image when the mouse
passed over it. To do them, a script reached into the page and changed it.

The page, as scripts see it, is the \textbf{DOM}, the Document Object Model: a
tree of objects, one per element, that the browser keeps in memory and redraws
whenever it changes. The W3C standardized it in 1998, but Netscape and Internet
Explorer each kept their own version for years, and a script that worked in one
often failed in the other. In January 2006, John Resig presented
\textbf{jQuery} at BarCamp NYC. It hid the differences between browsers behind
one short syntax: find these elements, change them. For the next ten years,
most of the web was written that way.
\end{history}

Here is a page written that way: a count of unread messages, a button to mark
them all as read, and a list of conversations with the unread ones in bold.

\begin{code}[language={},morestring={[b]"}]
<span id="badge">2</span>
<button id="read-all">mark all as read</button>
<ul id="threads">
  <li class="unread">Ana</li>
  <li class="unread">Bruno</li>
</ul>
\end{code}

And its script, in jQuery. \texttt{\$('\#badge')} finds the element whose
\texttt{id} is \texttt{badge}, \texttt{.text(...)} changes what it shows, and
\texttt{.on('click', ...)} registers a callback for clicks, the same kind of
callback as in Chapter~3:

\begin{code}
let unread = 2;

$('#threads li').on('click', function () {
  if (!$(this).hasClass('unread')) return;
  $(this).removeClass('unread');
  unread -= 1;
  $('#badge').text(unread);
});

$('#read-all').on('click', function () {
  $('#threads li').removeClass('unread');
  unread = 0;
});
\end{code}

The number of unread messages now lives in two places: in the variable
\texttt{unread}, which is what the program believes, and in the text of the
badge, which is what the user sees. Nothing ties them together. Each handler
has to remember every place that shows the number and change each one itself.
Figure~\ref{fig:dom-hands} shows what each handler touches.

\begin{figure}[htbp]
\centering
\begin{tikzpicture}[x=1cm, y=1cm]
  \tikzset{
    el/.style={hf/reader, font=\ttfamily\scriptsize, minimum height=6mm, inner sep=2pt},
    hd/.style={hf/bus, font=\sffamily\footnotesize, minimum width=30mm},
    touch/.style={hf/flow, black!55},
  }
  % the DOM
  \node[hf/note, font=\sffamily\footnotesize\bfseries, text=hfReader] at (0.4,1.3) {the page (the DOM)};
  \node[el] (doc) at (0.4,0.4) {page};
  \node[el] (badge) at (-2.4,-1.0) {span\#badge\\{\sffamily ``2''}};
  \node[el] (btn) at (0.0,-1.0) {button\#read-all};
  \node[el] (ul) at (2.6,-1.0) {ul\#threads};
  \node[el] (ana) at (1.6,-2.4) {li.unread\\{\sffamily Ana}};
  \node[el] (bruno) at (3.6,-2.4) {li.unread\\{\sffamily Bruno}};
  \draw[black!30] (doc) -- (badge); \draw[black!30] (doc) -- (btn);
  \draw[black!30] (doc) -- (ul);
  \draw[black!30] (ul) -- (ana); \draw[black!30] (ul) -- (bruno);
  % the program
  \node[hf/note, font=\sffamily\footnotesize\bfseries, text=hfStore] at (7.2,1.3) {the program};
  \node[hf/store, font=\ttfamily\scriptsize, minimum width=22mm] (var) at (7.2,-1.0) {let unread = 2};
  % the handlers
  \node[hd] (h1) at (-0.4,-4.4) {click on a name};
  \node[hd] (h2) at (4.4,-4.4) {mark all as read};
  \draw[touch] (h1) -- (badge);
  \draw[touch] (h1) -- (ana);
  \draw[touch] (h1.east) to[out=10, in=-120] (var.south west);
  \draw[touch] (h2) -- (ana);
  \draw[touch] (h2) -- (bruno);
  \draw[touch] (h2) -- (var);
  \draw[hf/pain, dashed] (h2.south) -- ++(0,-0.5) -| (-3.4,-1.0) -- (badge.west);
  \node[text=hfAlert, font=\sffamily\bfseries] at (-3.4,-3.0) {$\times$};
  \node[hf/note, text=hfAlert, anchor=north] at (0.6,-5.3) {nobody changes the badge here};
\end{tikzpicture}
\caption{What each handler of the script touches. The page is a tree of
objects, and each handler reaches into it by hand. ``Click on a name'' changes
the name, the variable and the badge. ``Mark all as read'' changes the names and
the variable, and forgets the badge.}
\label{fig:dom-hands}
\end{figure}

Click Ana, then the button, and the page ends up disagreeing with the program
(Figure~\ref{fig:dom-moments}).

\begin{figure}[htbp]
\centering
\begin{tikzpicture}[x=1cm, y=1cm]
  \tikzset{
    when/.style={hf/note, anchor=west, align=left, text width=36mm},
    head/.style={hf/note, font=\sffamily\footnotesize\bfseries, align=center},
    val/.style={font=\ttfamily\small},
  }
  \node[head, text=hfStore] at (5.6,0.9) {in the program\\\texttt{unread}};
  \node[head, text=hfReader] at (8.4,0.9) {on the page\\the badge};
  \node[head, text=hfReader] at (11.4,0.9) {on the page\\names in bold};
  \draw[black!12] (-0.1,0.35) -- (13.2,0.35);
  \foreach \r in {1,2} \draw[black!12] (-0.1,{-1.0*\r+0.45}) -- (13.2,{-1.0*\r+0.45});
  \node[when] at (0,0) {\textbf{1.} the page loads};
  \node[val] at (5.6,0) {2}; \node[val] at (8.4,0) {2};
  \node[val, font=\sffamily\small] at (11.4,0) {\textbf{Ana}, \textbf{Bruno}};
  \node[when] at (0,-1) {\textbf{2.} click on Ana};
  \node[val] at (5.6,-1) {1}; \node[val] at (8.4,-1) {1};
  \node[val, font=\sffamily\small] at (11.4,-1) {Ana, \textbf{Bruno}};
  \node[when] at (0,-2) {\textbf{3.} click on ``mark all as read''};
  \node[val] at (5.6,-2) {0};
  \node[val, text=hfAlert, font=\ttfamily\small\bfseries] at (8.4,-2) {1};
  \node[val, font=\sffamily\small] at (11.4,-2) {Ana, Bruno};
  \node[hf/note, text=hfAlert] at (8.4,-2.55) {the program says 0};
\end{tikzpicture}
\caption{The page from Figure~\ref{fig:dom-hands}, one click at a time. After
the third row, the badge says there is one unread message, and the program,
the list and the user all know there are none.}
\label{fig:dom-moments}
\end{figure}

That's the whole bug: one missing line, \texttt{\$('\#badge').text(0)}, in the
second handler. Here it's easy to see, with two handlers and two places to
update. A real page has dozens of handlers and the same number shown in
several places, and every new feature adds a handler that has to know all of
them. Sooner or later someone forgets one, and the screen starts showing
something the program no longer believes. jQuery made changing the page easy,
but it didn't change who had to remember: the programmer, by hand, in every
handler.

\begin{history}{2014 \textperiodcentered{} The message that wasn't there}
At F8, Facebook's yearly conference for developers, in April 2014, Facebook
engineers Tom Occhino and Jing Chen explained
why they had stopped building their pages that way. Their example was a bug
that kept coming back: the messages icon showed a count of unread messages, the
user clicked it, and there were no unread messages. The count, the list of
conversations and the open chat windows each kept their own idea of what had
been read, and each event had to update all of them. Engineers fixed the bug,
it stayed fixed for a while, and it returned with the next feature that touched
any of those places.
\end{history}

\begin{figure}[htbp]
\centering
\begin{tikzpicture}
  % by hand
  \node[hf/note, font=\sffamily\footnotesize\bfseries] at (0,1.2) {by hand};
  \node[hf/source, minimum width=30mm] (ev1) at (0,0) {a message is read};
  \node[hf/reader, minimum width=20mm] (badge) at (-2.3,-2.2) {unread count};
  \node[hf/reader, minimum width=20mm] (list) at (0,-2.2) {thread list};
  \node[hf/reader, minimum width=20mm] (chat) at (2.3,-2.2) {chat window};
  \draw[hf/ask] (ev1) -- (list);
  \draw[hf/ask] (ev1) -- (chat);
  \draw[hf/pain, dashed] (ev1) -- (badge) node[pos=0.55, text=hfAlert, font=\sffamily\bfseries] {$\times$};
  \node[hf/note, text=hfAlert] at (-2.3,-3.2) {forgotten: shows 1};
  % from state
  \begin{scope}[xshift=7.6cm]
  \node[hf/note, font=\sffamily\footnotesize\bfseries] at (0,1.2) {from the state};
  \node[hf/source, minimum width=30mm] (ev2) at (0,0) {a message is read};
  \node[hf/store, minimum width=22mm] (st) at (0,-1.5) {state};
  \node[hf/reader, minimum width=56mm] (scr) at (0,-3.0) {count, list and window, drawn again};
  \draw[hf/save] (ev2) -- node[hf/note, right] {change it once} (st);
  \draw[hf/flow, draw=hfStore] (st) -- node[hf/note, right] {draw from it} (scr);
  \end{scope}
\end{tikzpicture}
\caption{Left: each event updates every place that shows the data, and a
forgotten place keeps showing an old value. Right: the event changes the state
once, and everything on the screen is drawn from the state again.}
\end{figure}

\section{The screen as a function of the state}
\label{sec:click-models}

The fix came from an older way of building pages, one in which a count of
unread messages that nobody updated couldn't happen, because nothing on the
page lived long enough to go stale.

\begin{history}{1993--2005 \textperiodcentered{} Pages built on the server}
The first web servers only sent files. In 1993, the NCSA web server added CGI,
the Common Gateway Interface: for some addresses, instead of sending a file, the
server ran a program and sent back whatever the program printed. A page could
now be built from data at the moment someone asked for it. On June 8, 1995,
Rasmus Lerdorf announced the Personal Home Page Tools, a few such programs he
had written to count the visits to his online résumé. They grew into
\textbf{PHP}, a language written inside the HTML: the page is mostly HTML, and
wherever a value should appear, a short piece of code fetches it and prints it.
PHP became the most common way to build a page on the server, and in 2004 Mark
Zuckerberg wrote Facebook in it.

In that model the page is built from scratch on every request and thrown away
once it's sent. When a message is read, the next request builds a new page, and
the count, the list and everything else come from the same data at the same
moment. They can't disagree. The price is that every change, however small,
means asking the server for a whole new page and drawing it again. Ajax, from
Chapter~4, ended that around 2005: the page stayed open and changed itself in
place, which was much faster for the user and brought in the forgotten badge
of the previous section.
\end{history}

\begin{figure}[htbp]
\centering
\begin{tikzpicture}[x=1cm, y=1cm]
  \tikzset{
    st/.style={hf/base, font=\sffamily\scriptsize, minimum height=7mm, inner sep=2pt},
    ev/.style={st, draw=hfSource, fill=hfSource!10},
    sv/.style={st, draw=hfBus, fill=hfBus!10},
    pg/.style={st, draw=hfReader, fill=hfReader!10},
    sa/.style={st, draw=hfStore, fill=hfStore!10},
    yr/.style={hf/note, font=\sffamily\footnotesize\bfseries, anchor=west, align=left, text width=24mm},
    good/.style={hf/note, anchor=west, align=left, text=hfStore},
    bad/.style={hf/note, anchor=west, align=left, text=hfAlert},
    fl/.style={hf/flow, black!55},
  }
  % 1995
  \node[yr] at (-0.2,0) {1995\\[-1pt]{\mdseries server builds the page}};
  \node[ev] (c1) at (3.0,0) {click};
  \node[sv] (s1) at (5.5,0) {the server builds\\a whole new page};
  \node[pg] (p1) at (8.3,0) {the browser draws\\it from scratch};
  \draw[fl] (c1) -- (s1);
  \draw[fl] (s1) -- (p1);
  \node[good] at (10.0,0.2) {nothing to forget};
  \node[bad] at (10.0,-0.2) {slow: a page per click};
  \draw[black!12] (-0.2,-0.8) -- (13.8,-0.8);
  % 2005
  \node[yr] at (-0.2,-1.95) {2005\\[-1pt]{\mdseries the page changes itself}};
  \node[ev] (c2) at (3.0,-1.95) {click};
  \node[sv] (h2) at (5.0,-1.95) {handler};
  \node[sa, minimum height=5mm, minimum width=15mm] (v2) at (7.6,-1.35) {variable};
  \node[pg, minimum height=5mm, minimum width=15mm] (a2) at (7.6,-1.95) {list};
  \node[pg, minimum height=5mm, minimum width=15mm] (b2) at (7.6,-2.55) {badge};
  \draw[fl] (c2) -- (h2);
  \draw[fl] (h2.east) -- (v2.west);
  \draw[fl] (h2.east) -- (a2.west);
  \draw[hf/pain, dashed] (h2.east) -- (b2.west);
  \node[text=hfAlert, font=\sffamily\bfseries\scriptsize] at (6.45,-2.38) {$\times$};
  \node[good] at (10.0,-1.75) {fast};
  \node[bad] at (10.0,-2.15) {something to forget};
  \draw[black!12] (-0.2,-3.1) -- (13.8,-3.1);
  % 2013
  \node[yr] at (-0.2,-4.0) {2013\\[-1pt]{\mdseries the page is drawn from the state}};
  \node[ev] (c3) at (3.0,-4.0) {click};
  \node[sa] (s3) at (4.6,-4.0) {state};
  \node[sv] (d3) at (6.4,-4.0) {a new\\description};
  \node[pg] (p3) at (8.4,-4.0) {the page: only\\what changed};
  \draw[fl] (c3) -- (s3);
  \draw[fl] (s3) -- (d3);
  \draw[fl] (d3) -- (p3);
  \node[good] at (10.0,-3.8) {nothing to forget};
  \node[good] at (10.0,-4.2) {fast};
\end{tikzpicture}
\caption{Three ways to answer a click. On the server, every click builds a new
page from the data: nothing can go stale, but every click costs a whole page.
In the browser, by hand, every handler changes the page in place: fast, but each
one has to remember every spot. React keeps the first model's rule, the page is
always drawn from the data, and runs it in the browser.}
\label{fig:three-models}
\end{figure}

\begin{question}{Isn't that server-side rendering?}
Yes. Building the HTML on the server for each request is what
\emph{server-side rendering} means, but in 1995 nobody called it that: it was
simply how the web worked. The name became necessary only once pages started
being built in the browser, by JavaScript, in the early 2010s, and the old way
needed a word to tell it apart, the same way ``macrotask'' was coined in
Chapter~3. Today the term usually means something narrower: running the same
components on the server to send the first HTML, so the page appears before the
JavaScript arrives, and letting them take over in the browser after that.
\end{question}

\begin{history}{2010--2013 \textperiodcentered{} Jordan Walke}
Facebook's server code was still PHP. In February 2010 it released XHP, an
extension written by Marcel Laverdet that let PHP code contain HTML elements as
values, built from small pieces and escaped by the language itself. On the
server the rule was simple: every request built the whole page again from the
data, so the page could never disagree with it. In the browser, since Ajax, the
rule was the opposite: the page stayed open for hours, and every handler
patched it by hand. That's where the phantom badge came from.

Jordan Walke, an engineer at Facebook, wanted the server's rule in the browser,
without giving back what Ajax had won: on every change, describe the whole page
again from the data, the way PHP did on every request, but without asking the
server for anything. The obvious way to do it doesn't work. Throwing away the
page's elements and building them again is slow, and it loses everything the
browser keeps inside them: the text half typed into a field, the position of
the cursor, how far the page was scrolled. Walke rebuilt something cheaper
instead: a description of the page made of plain JavaScript objects. He
compared it with the previous description and changed in the real page only
what was different. That comparison, later called the \emph{virtual DOM}, is
what made ``build everything again on every change'' affordable.

He built a prototype called FaxJS, renamed it
FBolt inside Facebook in 2012, and then React. It ran parts of Facebook's news
feed and, in 2012, Instagram's web site. On May 29, 2013, Walke and Tom
Occhino released it at JSConf US. The first reactions were mostly hostile,
because React put HTML inside JavaScript, and separating the two was what
everyone had been taught.
\end{history}

\begin{concept}{UI as a function of state}
The program never changes the page. It keeps its data in a state, and it
describes what the page should look like for a given state. When the state
changes, the description is computed again, and the library makes the page
match it.

\textbf{What it fixes:} only one piece of code decides what each part of the
screen shows, so there's no second place to forget. A badge left showing a
message that was already read can't happen.

\textbf{What it costs:} the description is recomputed on every change, even
when almost nothing changed, and someone has to work out what to actually do to
the page. That someone is React.
\end{concept}

The rest of this chapter takes React apart in the order the code meets it: what
you write, what a tool turns it into, how it reaches the browser, and what
happens there when someone clicks.

\section{A component, line by line}

In React, the description is written as a \emph{component}: a function that
returns what to draw. This is a whole one, a button that counts its clicks:

\begin{code}
import { useState } from 'react';

export function Counter() {
  const [count, setCount] = useState(0);
  return (
    <button onClick={() => setCount(count + 1)}>
      clicked {count} times
    </button>
  );
}
\end{code}

Line~1 takes \texttt{useState} from the \texttt{react} package. It isn't an
object or another name for something: it's a function, one of the few that
React gives components to call, and it belongs to a family called
\emph{hooks}.

\begin{history}{2018--2019 \textperiodcentered{} Hooks}
Until 2019, a React component that had state had to be a class, with methods
that React called at fixed moments: when the component appeared, when it
updated, when it went away. Related logic ended up split across those methods,
and sharing it between components needed patterns that wrapped components
inside other components. At React Conf, in October 2018, Sophie Alpert and Dan
Abramov presented \textbf{Hooks}, and React 16.8 shipped them in February 2019:
functions like \texttt{useState} that a plain function component calls to get
state and other features.
\end{history}

Line~3 declares the component, an ordinary function whose name starts with a
capital letter. Nothing in the file calls it, and nothing in your project will:
React does. This is the inversion from Chapter~3. \texttt{Counter} is code you
plug into React, and React calls it whenever the button has to be drawn: once
when the page starts, and again every time its state changes. React calls
itself a library, but by the definition in Chapter~3 it's a framework.

Line~4 is where the state comes in, and it does three things at once.
\texttt{useState(0)} asks React for this component's state, with \texttt{0} as
the value to start with. React keeps the value, not your function:
\texttt{Counter} ends after every drawing, and its local variables end with it,
so the number has to live somewhere that survives between calls.
\texttt{useState} returns a list of two things, and the square brackets, a
feature of ES2015 called destructuring, take them out of the list and name
them. The names are yours; \texttt{[n, setN]} would work the same.
\texttt{count} is a plain number, the value for this drawing.
\texttt{setCount} is a function that asks React to keep a new value and to draw
the component again.

\texttt{setCount} doesn't change \texttt{count}. \texttt{count} is a
\texttt{const}, and it stays \texttt{0} until this run of \texttt{Counter} is
over. The new number reaches the screen through the next run, the one React
starts after \texttt{setCount}: there, \texttt{useState} returns the new value.

Lines~5 to 9 describe the button: an element with a click handler and a text
with the number in it. It looks like HTML, but it's JSX, the subject of the
next section, and what it produces when it runs isn't a button on the page. It's
a few plain JavaScript objects that describe one. The click handler is an arrow
function that calls \texttt{setCount} with the next number.

No line touches the page. Nothing here finds the button and changes its text,
the way the jQuery handlers did. That's React's job. How it does it, from the
moment the page arrives in the browser to the click, is
Section~\ref{sec:click}, once the next two sections have shown what the JSX
becomes and how the code gets to the browser.

Everything in the rest of this chapter that looks strange about React comes
from one fact: React, not you, decides when \texttt{Counter} runs.

\section{JSX}

The \texttt{<button>} inside the function above isn't JavaScript. It's
\textbf{JSX}, and no browser understands it.

JSX came from the same place as React. XHP's author, Marcel Laverdet, had
written a small project that added XML literals to JavaScript, and Walke forked
it into JSX, changing it into a shorter way of writing nested function calls.
Before running, every JSX element is turned into a call. This is
\texttt{Counter} after that step:

\begin{code}
import { useState } from "react";
import { jsxs as _jsxs } from "react/jsx-runtime";
export function Counter() {
  const [count, setCount] = useState(0);
  return /* @__PURE__ */ _jsxs("button", {
    onClick: () => setCount(count + 1),
    children: [
      "clicked ",
      count,
      " times"
    ]
  });
}
\end{code}

Two programs act on this code, at two different moments, and neither of them
draws anything.

The first is the build tool, the subject of the next section. It works before
the browser ever gets the file, and it works on text: it reads the JSX and
writes the JavaScript above, without running any of it. The element became a
call to \texttt{\_jsxs} with the tag name and an object of properties, and its
contents became a list of children. Everything between braces in JSX is
ordinary JavaScript, passed along as a value. (The \texttt{@\_\_PURE\_\_}
comment tells the bundler the call has no side effects, so it can be dropped if
nothing uses its result.)

The second moment is in the browser, each time React calls \texttt{Counter}.
What runs there is the code above, plain JavaScript, and the browser's engine
runs it like any other script. \texttt{\_jsxs} is an ordinary function exported
by React, and it doesn't draw either: it builds an object and returns it.
Called with the same arguments in Node, with \texttt{count} at 0, it returns
this:

\begin{terminal}[]
{
  '$$typeof': Symbol(react.transitional.element),
  type: 'button',
  key: null,
  props: {
    onClick: [Function: onClick],
    children: [ 'clicked ', 0, ' times' ]
  },
  _owner: null,
  _store: {}
}
\end{terminal}

This object is the description from Section~\ref{sec:click-models}, and React
calls it an \emph{element}. \texttt{type} says which tag, \texttt{props} holds
the handler and the children, and \texttt{\$\$typeof} marks the object as made
by React. Nothing in it is graphical. It's data, as plain as the JSON a server
sends.

The button on the screen comes from a third step, which belongs to
\texttt{react-dom}, the part of React that knows about browsers. It reads the
element and makes the same DOM calls the jQuery handlers made by hand in
Section~\ref{sec:by-hand}: it creates the button and its texts and puts them in
the page. Figure~\ref{fig:jsx-steps} shows the three steps, with the DOM calls
as they were logged in a simulated browser. \texttt{onClick} isn't in any of
those calls. Section~\ref{sec:click} shows where it goes.

\begin{figure}[htbp]
\centering
\begin{tikzpicture}[x=1cm, y=1cm]
  \tikzset{
    who/.style={hf/note, anchor=west, align=left, text width=29mm},
    cd/.style={hf/base, draw=black!40, fill=black!4, font=\ttfamily\scriptsize, inner sep=2pt, align=left},
    pg/.style={hf/reader, font=\ttfamily\scriptsize, minimum width=30mm, minimum height=6mm, inner sep=2pt},
    head/.style={hf/note, font=\sffamily\footnotesize\bfseries, align=center},
    say/.style={hf/note, font=\sffamily\scriptsize, text=hfBus, text width=24mm, above},
  }
  \node[head] at (1.45,1.2) {who, and when};
  \node[head] at (5.0,1.2) {gets};
  \node[head] at (11.2,1.2) {gives};
  \draw[black!12] (-0.1,0.8) -- (13.9,0.8);
  % 1
  \node[who] at (0,0) {\textbf{1.} The build tool, before the browser gets the file.};
  \node[cd] (g1) at (5.0,0) {<button onClick=\{...\}>\\\ \ clicked \{count\} times\\</button>};
  \node[cd] (o1) at (11.2,0) {\_jsxs("button", \{\\\ \ onClick: ...,\\\ \ children: [...]\\\})};
  \draw[hf/flow, draw=hfBus] (g1) -- node[say] {rewrites text,\\runs nothing} (o1);
  \draw[black!12] (-0.1,-1.05) -- (13.9,-1.05);
  % 2
  \node[who] at (0,-1.9) {\textbf{2.} \texttt{\_jsxs}, in the browser, each time React calls \texttt{Counter}.};
  \node[cd] (g2) at (5.0,-1.9) {"button",\\\{ onClick,\\\ \ children \}};
  \node[cd] (o2) at (11.2,-1.9) {\{ type: "button",\\\ \ props: \{ onClick,\\\ \ \ \ children \} \}};
  \draw[hf/flow, draw=hfBus] (g2) -- node[say] {builds an object,\\draws nothing} (o2);
  \draw[black!12] (-0.1,-2.85) -- (13.9,-2.85);
  % 3
  \node[who] at (0,-4.3) {\textbf{3.} \texttt{react-dom}, in the browser, after \texttt{Counter} returns.};
  \node[cd] (g3) at (4.6,-4.3) {\{ type: "button",\\\ \ props: \{ ... \} \}};
  \node[cd] (o3) at (11.2,-4.3) {document.createTextNode("clicked ")\\document.createTextNode("0")\\document.createTextNode(" times")\\document.createElement("button")\\button.appendChild(\#text)\\button.appendChild(\#text)\\button.appendChild(\#text)\\div.appendChild(button)};
  \draw[hf/flow, draw=hfBus] (g3) -- node[say] {calls the DOM} (o3);
  \node[pg] (b) at (11.2,-6.25) {clicked 0 times};
  \draw[hf/flow, draw=hfReader] (o3) -- (b);
  \node[hf/note, font=\sffamily\scriptsize, text=hfReader, anchor=east] at ([xshift=-2mm]b.west) {on the screen};
\end{tikzpicture}
\caption{From JSX to a button. The build tool rewrites text before the browser
gets it. In the browser, \texttt{\_jsxs} runs each time React calls
\texttt{Counter} and returns an object. \texttt{react-dom} reads the object and
calls the DOM. Only the last step touches the page.}
\label{fig:jsx-steps}
\end{figure}

That's also why JSX writes \texttt{className} where HTML writes
\texttt{class}: the attributes become properties of a JavaScript object, and
\texttt{className} is what the DOM calls that property, because \texttt{class}
was a reserved word in early JavaScript.

The objection from 2013, that HTML and JavaScript should be kept apart, was
wrong. Both of them describe the same screen, and splitting them by language
puts the two halves of every component in different files. JSX keeps together
the things that change together. Its real cost is the one the browser imposes:
since no browser runs JSX, a React project can't run without a build step.

\section{Why a front end has a build}

\begin{history}{2009--2026 \textperiodcentered{} Bundlers}
For most of its life, the browser had no modules. A page loaded its scripts
with one \texttt{<script>} tag each, in the right order, and they talked through
global variables. Once Node gave JavaScript \texttt{require}, people wanted it in
the browser too. \textbf{Browserify} (2011) followed every \texttt{require} in a
project and wrote all the files into one, a \emph{bundle}. \textbf{webpack},
which Tobias Koppers started in 2012 and released as 1.0 in 2014, did the same
for every kind of file: JavaScript, styles, images. Babel (2014, first called
6to5) translated new syntax, and later JSX, into something every browser could
run. In 2016, Facebook released Create React App, which set all of it up with
one command. That's how React projects started for most of a decade.

Browsers learned ES modules between 2017 and 2018: Safari first, then Chrome
61 in September 2017, then Firefox 60 in May 2018. In April 2020, Evan You,
who had created the Vue framework in 2014, released \textbf{Vite}. In
February 2025, the React team deprecated Create React App, which had no active
maintainers left, and pointed to frameworks and to build tools like Vite
instead. Vite 8, in March 2026, replaced its bundler with Rolldown, written in
Rust.
\end{history}

webpack's dev server bundled the whole project before showing the first page,
and again, at least in part, after every change. On a large project that took
long enough to go get a coffee. Vite's idea was that once browsers understood
\texttt{import}, the dev server didn't need to bundle anything. The browser
asks for \texttt{main.jsx}, finds an \texttt{import} of \texttt{Counter.jsx} in
it and asks for that, and the server transforms each file at the moment it's
requested. A project with a thousand files and a page that uses forty costs
forty transforms.

\begin{figure}[htbp]
\centering
\begin{tikzpicture}[x=1mm, y=1mm]
  % bundle first
  \node[hf/note, anchor=west, font=\sffamily\footnotesize\bfseries] at (0,40) {bundle first};
  \foreach \i in {0,...,7} \node[hf/source, minimum width=7mm, minimum height=5mm, inner sep=1pt, font=\sffamily\tiny] (a\i) at (5+\i*8,30) {};
  \node[hf/bus, minimum width=22mm] (b) at (90,30) {one bundle};
  \node[hf/reader, minimum width=18mm] (br1) at (130,30) {browser};
  \draw[hf/flow, black!60] (65,30) -- (b);
  \draw[hf/flow, black!60] (b) -- (br1);
  \node[hf/note] at (34,23) {every file, before the first page};
  % on demand
  \node[hf/note, anchor=west, font=\sffamily\footnotesize\bfseries] at (0,12) {on demand};
  \foreach \i in {0,...,7} \node[hf/source, minimum width=7mm, minimum height=5mm, inner sep=1pt, opacity=0.35] at (5+\i*8,2) {};
  \node[hf/source, minimum width=7mm, minimum height=5mm, inner sep=1pt] (m) at (5,2) {};
  \node[hf/source, minimum width=7mm, minimum height=5mm, inner sep=1pt] (c) at (21,2) {};
  \node[hf/bus, minimum width=22mm] (srv) at (90,2) {dev server};
  \node[hf/reader, minimum width=18mm] (br2) at (130,2) {browser};
  \draw[hf/ask] ([yshift=1.5mm]br2.west) -- ([yshift=1.5mm]srv.east);
  \draw[hf/answer] ([yshift=-1.5mm]srv.east) -- ([yshift=-1.5mm]br2.west);
  \node[hf/note, anchor=south] at (110,5) {asks per file};
  \draw[hf/flow, black!60] (m.south) |- ++(0,-5) -| ([xshift=-4mm]srv.south);
  \draw[hf/flow, black!60] (c.south) |- ++(0,-3) -| ([xshift=-2mm]srv.south);
  \node[hf/note] at (55,-9) {only the files the page imports, as it imports them};
\end{tikzpicture}
\caption{Two dev servers. Above, the classic one builds the whole project into
one file first. Below, Vite serves source files as ES modules and transforms
each one when the browser asks for it.}
\end{figure}

The dev server still does a lot of work you don't see. When the browser asks
for \texttt{Counter.jsx}, it gets plain JavaScript: the JSX has become function
calls, as above, and the import of \texttt{react} now points to a path,
\url{/node_modules/.vite/deps/react.js}. A browser
can't look up a bare name like \texttt{react} the way Node can, so the server
rewrites it. And React itself is still published as CommonJS, so Vite converts
it to an ES module once and keeps the result in that folder.

In production, Vite bundles anyway. A browser that discovers imports one file
at a time pays a round trip to the server for each level of the import graph,
and over a real network that's slower than downloading one file. So the dev
server and the production build work in opposite ways, and a project can run
fine in one and break in the other. The build of the counter above turned 15
modules into a single file of 219.58~kB, 68.64~kB compressed. Almost all of it
is React.

The \texttt{<script>} tag era of React is over too. React 19, in December 2024,
stopped publishing the build that defined a global \texttt{React} variable for
pages without a build step. Vite is the right default for a React project today.
Create React App is dead and nobody should start a project with it.

\section{From the server to a click}
\label{sec:click}

All the pieces are on the table: a component, JSX turned into function calls, a
server that delivers them. What's missing is the order in which they meet in
the browser. Besides \texttt{Counter.jsx}, the counter needs two more files. One
is the page itself, almost empty:

\begin{code}[language={},morestring={[b]"}]
<!doctype html>
<html>
  <body>
    <div id="root"></div>
    <script type="module" src="/src/main.jsx"></script>
  </body>
</html>
\end{code}

The other is the entry point, \texttt{main.jsx}, the one place where your code
calls React instead of the other way around:

\begin{code}
import { createRoot } from 'react-dom/client';
import { Counter } from './Counter.jsx';

createRoot(document.getElementById('root')).render(<Counter />);
\end{code}

\texttt{createRoot} takes the empty \texttt{div} and hands it to React.
\texttt{render} gives React the first description, \texttt{<Counter />}, which
says ``a \texttt{Counter} goes here'' without calling it. From then on, React
is in charge of everything inside that \texttt{div}.

Those are the three files you write, and none of them reaches the browser as
you wrote it. What arrives depends on which of the two roads from the previous
section you're on: the dev server while you work, or the build when the project
goes to users. Figure~\ref{fig:two-roads} puts them side by side, with the
requests a browser makes for the counter on each.

\begin{figure}[htbp]
\centering
\begin{tikzpicture}[x=1cm, y=1cm]
  \tikzset{
    req/.style={anchor=west, font=\ttfamily\scriptsize, inner sep=1pt},
    ans/.style={anchor=north west, hf/note, font=\sffamily\scriptsize, inner sep=1pt, align=left},
    tot/.style={hf/note, font=\sffamily\footnotesize\bfseries},
  }
  \draw[black!15] (7,0.9) -- (7,-6.3);
  % dev
  \node[hf/bus, minimum width=55mm] (dev) at (3.4,0.3) {in dev: the Vite dev server\\{\scriptsize\ttfamily npm run dev}};
  \foreach \y/\r/\a in {
      0/{GET /}/{your index.html, plus two scripts Vite added},
      1/{GET /@vite/client, /@react-refresh}/{Vite's code for updating the page on save},
      2/{GET /src/main.jsx}/{main.jsx, transformed when asked},
      3/{GET .../.vite/deps/react-dom\_client.js}/{React's development build},
      4/{GET /src/Counter.jsx}/{Counter.jsx, transformed when asked},
      5/{\ldots{} and five more}/{one request per import}} {
    \node[req] at (0.4,{-0.75-0.72*\y}) {\r};
    \node[ans] at (0.75,{-0.88-0.72*\y}) {\a};
  }
  \node[tot] (tdev) at (3.4,-5.35) {11 requests, about 3.6 MB};
  % production
  \node[hf/bus, minimum width=55mm] (prod) at (10.6,0.3) {in production: any web server\\{\scriptsize\ttfamily npm run build}{\scriptsize, once, before}};
  \foreach \y/\r/\a in {
      0/{GET /}/{dist/index.html: the div and one script tag},
      1/{GET /assets/index-BgWl9ddg.js}/{React and your two files, already\\bundled, JSX already turned into calls}} {
    \node[req] at (7.6,{-0.75-0.72*\y}) {\r};
    \node[ans] at (7.95,{-0.88-0.72*\y}) {\a};
  }
  \node[ans, anchor=north west, text width=58mm] at (7.6,-2.75) {The server transforms nothing. It sends the
    files the build wrote, and any server that sends files will do.};
  \node[tot] (tprod) at (10.6,-5.35) {2 requests, 219.58~kB of script};
  % both
  \node[hf/reader, minimum width=60mm] (run) at (7,-6.9) {the browser runs \texttt{main.jsx}};
  \draw[hf/flow, draw=hfReader] (tdev.south) |- (run.west);
  \draw[hf/flow, draw=hfReader] (tprod.south) |- (run.east);
\end{tikzpicture}
\caption{The two roads from your files to the browser, measured on the counter.
In dev, the page arrives almost as you wrote it, and the browser follows the
imports one file at a time while Vite transforms each. In production, the build
did that work earlier and left two files. Both roads end with the browser
running \texttt{main.jsx}.}
\label{fig:two-roads}
\end{figure}

In dev, Vite adds two scripts to the top of your page. \texttt{/@vite/client}
opens a WebSocket to the dev server, so that when you save a file the server
can tell the page to load the new version, and \texttt{/@react-refresh} lets
React swap a component without losing its state. Then the browser reads
\texttt{<script src="/src/main.jsx">}, asks for it, finds its imports and asks
for those, one level at a time. Most of the 3.6~MB is React's development
build, which is not minified and is full of checks that print warnings. The
JSX becomes calls to \texttt{\_jsxDEV} instead of \texttt{\_jsxs}: the same
call, with the file name and line number added, so those warnings can point at
your code.

In production, the work happens before anyone asks. \texttt{vite build} reads
\texttt{index.html}, follows every import from \texttt{main.jsx}, and writes
two files to a folder called \texttt{dist}. This is the page:

\begin{code}[language={},morestring={[b]"}]
<!doctype html>
<html>
  <script type="module" crossorigin src="/assets/index-BgWl9ddg.js"></script>

  <body>
    <div id="root"></div>
  </body>
</html>
\end{code}

The other file is the script, with React's production build and your two files
inside, the JSX already turned into function calls. The letters in its name
come from its contents, so a new build gets a new name, and a browser holding
the old one in its cache can't run stale code.

On both roads, the \texttt{div} arrives empty: no button, no number, no
listener. The page doesn't arrive with the state in it. The button exists only
after the script runs, and from the moment \texttt{main.jsx} runs the two roads
are the same. Figure~\ref{fig:click-chain} follows the counter from there to
the first drawing.

\begin{figure}[htbp]
\centering
\begin{tikzpicture}[x=1cm, y=1cm]
  \tikzset{
    pg/.style={hf/reader, font=\ttfamily\scriptsize, minimum width=25mm, minimum height=5mm, inner sep=1.5pt, align=center},
    cd/.style={hf/base, draw=black!40, fill=black!4, font=\ttfamily\scriptsize, minimum width=26mm, minimum height=5mm, inner sep=1.5pt, align=center},
    kp/.style={hf/store, font=\ttfamily\scriptsize, minimum width=15mm, minimum height=5mm, inner sep=1.5pt},
    rc/.style={hf/note, font=\sffamily\scriptsize, text=hfBus, align=center, text width=27mm},
    when/.style={hf/note, anchor=west, align=left, text width=31mm},
    head/.style={hf/note, font=\sffamily\footnotesize\bfseries, align=center},
    none/.style={hf/note, text=black!35, font=\sffamily\scriptsize},
    hot/.style={draw=hfAlert, line width=1pt},
  }
  \def\rh{1.8}
  \node[head, text=hfReader] at (4.65,0.95) {the page (DOM)};
  \node[head, text=hfBus] at (7.6,0.95) {React};
  \node[head] at (10.4,0.95) {your code};
  \node[head, text=hfStore] at (12.9,0.95) {kept by\\React};
  \draw[black!12] (-0.1,0.55) -- (13.9,0.55);
  \foreach \r in {1,...,3} \draw[black!12] (-0.1,{-\rh*\r+0.9}) -- (13.9,{-\rh*\r+0.9});
  % 1
  \node[when] at (0,0) {\textbf{1.} The HTML arrives.};
  \node[pg] at (4.65,0) {div\#root\\{\sffamily (empty)}};
  \node[none] at (7.6,0) {not loaded yet};
  \node[none] at (10.4,0) {not loaded yet};
  \node[none] at (12.9,0) {---};
  % 2
  \node[when] at (0,{-1*\rh}) {\textbf{2.} The script arrives and \texttt{main.jsx} runs.};
  \node[pg] at (4.65,{-1*\rh}) {div\#root\\{\sffamily (empty)}};
  \node[rc] at (7.6,{-1*\rh}) {loaded; receives the first description};
  \node[cd] at (10.4,{-1*\rh}) {main.jsx:\\render(<Counter />)};
  \node[none] at (12.9,{-1*\rh}) {---};
  % 3
  \node[when] at (0,{-2*\rh}) {\textbf{3.} React calls \texttt{Counter}; \texttt{useState} returns 0.};
  \node[pg] at (4.65,{-2*\rh}) {div\#root\\{\sffamily (empty)}};
  \node[rc] at (7.6,{-2*\rh}) {calls \texttt{Counter}};
  \node[cd] at (10.4,{-2*\rh}) {count = 0\\returns a description};
  \node[kp] at (12.9,{-2*\rh}) {count: 0};
  % 4
  \node[when] at (0,{-3*\rh}) {\textbf{4.} React builds the page from the description.};
  \node[pg] at (4.65,{-3*\rh}) {<button>\\clicked 0 times};
  \node[rc] at (7.6,{-3*\rh}) {creates the button and its text, keeps the description};
  \node[none] at (10.4,{-3*\rh}) {finished};
  \node[kp] at (12.9,{-3*\rh}) {count: 0};
\end{tikzpicture}
\caption{The counter from the HTML to the first drawing, the same on both roads.
The page arrives empty. Your code runs in rows 2 and 3, and never touches the page:
\texttt{main.jsx} hands React the first description, and \texttt{Counter}
returns the description of the button. The number lives in React (green), and
only React builds the page (blue).}
\label{fig:click-chain}
\end{figure}

After row~4, the browser's inspector shows the page as
\texttt{<div id="root"><button>clicked 0 times</button></div>}, and that isn't
all of it. React also attached things that never appear in the HTML
(Figure~\ref{fig:react-left}).

\begin{figure}[htbp]
\centering
\begin{tikzpicture}[x=1cm, y=1cm]
  \tikzset{
    el/.style={hf/reader, font=\ttfamily\scriptsize, minimum width=26mm, minimum height=6mm, align=center},
    hid/.style={hf/note, font=\sffamily\scriptsize, text=hfBus, align=left, anchor=west},
    rec/.style={hf/store, font=\sffamily\scriptsize, minimum width=30mm, align=center},
    link/.style={hf/flow, draw=hfBus, dashed},
  }
  \node[hf/note, font=\sffamily\footnotesize\bfseries, text=hfReader] at (0,1.7) {what the HTML shows};
  \node[hf/note, font=\sffamily\footnotesize\bfseries, text=hfBus] at (5.2,1.7) {what React attached};
  \node[hf/note, font=\sffamily\footnotesize\bfseries, text=hfStore] at (10.6,1.7) {what React keeps};
  \node[el] (root) at (0,0.4) {div\#root};
  \node[el] (btn) at (0,-1.6) {button\\{\sffamily clicked 0 times}};
  \draw[black!35] (root) -- (btn);
  \node[hid] (l1) at (2.2,0.4) {87 listeners, one per kind of event:\\\texttt{click}, \texttt{keydown}, \texttt{input}, \ldots};
  \draw[link] (l1.west) -- (root.east);
  \node[hid] (l2) at (2.2,-1.25) {no listener of its own};
  \node[hid] (l3) at (2.2,-1.75) {hidden: its props, \texttt{onClick} among them};
  \node[hid] (l4) at (2.2,-2.25) {hidden: a link to its record};
  \draw[link] (l2.west) -- (btn.east);
  \node[rec] (rb) at (10.6,-2.25) {record of the button};
  \node[rec] (rc) at (10.6,0.1) {record of \texttt{Counter}:\\the function \texttt{Counter},\\its hooks: \texttt{[0]}};
  \draw[link] (l4.east) -- (rb.west);
  \draw[hf/flow, draw=hfStore] (rb) -- node[hf/note, right] {parent} (rc);
\end{tikzpicture}
\caption{The page after the first drawing. The HTML shows only the
\texttt{div} and the button. React put the listeners on the \texttt{div}, not on
the button, and left two hidden properties on the button: its props, and a link
to an internal record. Each place in the tree has a record, and the record of
\texttt{Counter} keeps the function and its state.}
\label{fig:react-left}
\end{figure}

All of it comes from the same run. In the built counter, \texttt{createRoot}
registered 87 kinds of event on the \texttt{div}, \texttt{click} among them, and
none on the button. The button carries two extra properties whose names start
with \texttt{\_\_reactProps} and \texttt{\_\_reactFiber}. The first holds the
props, \texttt{onClick} and \texttt{children}. The second points to the
button's record, which React calls a \emph{fiber}, the name of the engine React
was rewritten around in React~16, in 2017. Going up from the button's record
leads to the record of \texttt{Counter}, which holds the \texttt{0}.

Now the click. Figure~\ref{fig:react-change} follows it from the user's finger
to the new text.

\begin{figure}[htbp]
\centering
\begin{tikzpicture}[x=1cm, y=-0.62cm]
  \tikzset{
    lane/.style={hf/base, minimum width=24mm, minimum height=9mm, font=\sffamily\footnotesize},
    call/.style={hf/flow, black!70},
    back/.style={hf/flow, black!45, dashed},
    lbl/.style={hf/note, above, font=\sffamily\scriptsize},
    mine/.style={draw=hfReader, fill=hfReader!25, line width=0.6pt},
    rnote/.style={hf/note, font=\sffamily\scriptsize, text=hfBus, align=left, anchor=west, fill=white, inner sep=1pt},
  }
  \def\U{0} \def\B{3.0} \def\R{6.6} \def\C{10.6}
  \node[lane, draw=hfSource, fill=hfSource!10] at (\U,0) {the user};
  \node[lane, draw=hfReader, fill=hfReader!10] at (\B,0) {the browser\\{\scriptsize and the page}};
  \node[lane, draw=hfBus, fill=hfBus!10] at (\R,0) {React};
  \node[lane, draw=hfReader, line width=1.2pt, fill=hfReader!25] at (\C,0) {your code};
  \foreach \x in {\U,\B,\R,\C} \draw[hf/lifeline] (\x,0.8) -- (\x,17.6);
  % your code runs only inside these two bars
  \draw[mine] (\C-0.09,4.2) rectangle (\C+0.09,7.0);
  \draw[mine] (\C-0.09,10.2) rectangle (\C+0.09,13.6);
  \node[hf/note, anchor=west, align=left, text=hfReader, font=\sffamily\scriptsize\bfseries] at (\C+0.2,5.6) {the click\\handler};
  \node[hf/note, anchor=west, align=left, text=hfReader, font=\sffamily\scriptsize\bfseries] at (\C+0.2,11.9) {\texttt{Counter}};
  % the click
  \draw[call] (\U,1.6) -- node[lbl] {clicks the button} (\B,1.6);
  \draw[call] (\B,2.8) -- node[lbl] {calls React's listener} (\R,2.8);
  \draw[call] (\R,4.2) -- node[lbl] {calls your \texttt{onClick}} (\C,4.2);
  \node[rnote, anchor=east, align=right] at (\R-0.12,3.75) {finds the button's\\hidden props};
  \draw[call] (\C,5.3) -- node[lbl] {\texttt{setCount(1)}} (\R,5.3);
  \node[rnote, anchor=east, align=right] at (\R-0.12,6.0) {keeps \texttt{count = 1},\\plans a new drawing};
  \draw[back] (\C,7.0) -- node[lbl] {the handler returns} (\R,7.0);
  \draw[back] (\R,8.0) -- node[lbl] {the click is over} (\B,8.0);
  \node[hf/note, font=\sffamily\scriptsize, text=hfAlert, fill=white] at (\B,9.0) {still shows 0};
  % the new drawing
  \draw[call] (\R,10.2) -- node[lbl] {calls \texttt{Counter()}} (\C,10.2);
  \draw[call] (\C,11.2) -- node[lbl] {\texttt{useState(0)}} (\R,11.2);
  \draw[back] (\R,12.2) -- node[lbl] {\texttt{[1, setCount]}} (\C,12.2);
  \draw[back] (\C,13.6) -- node[lbl] {returns a description} (\R,13.6);
  \node[rnote] at (\R+0.12,14.5) {compares it with the previous one:\\only the number differs};
  \draw[call] (\R,15.8) -- node[lbl] {changes one text: 0 $\to$ 1} (\B,15.8);
  \draw[call] (\B,17.0) -- node[lbl] {draws the screen} (\U,17.0);
\end{tikzpicture}
\caption{The click: who calls whom. Solid arrows are calls, dashed arrows are
returns. Your code runs twice, in the two blue bars, and both times because
React called it: first the click handler, which only asks React for a new
value, and then \texttt{Counter}, after the click is over. Neither touches the
page. Only React does, at the end.}
\label{fig:react-change}
\end{figure}

Traced in the simulated browser, the order is exactly that: the handler starts
and ends, the click event finishes, \texttt{Counter} runs and gets \texttt{1}
from \texttt{useState}, and only then does the page change. It changes once:
one piece of text, \texttt{"0"} becoming \texttt{"1"}, inside the same button
element, which React kept instead of building a new one.

\begin{question}{How does React know it has to call Counter, of all functions?}
Not from the click. The click only leads React from the clicked element to its
hidden props, and from the props to the \texttt{onClick} to call. The redraw
comes from \texttt{setCount}, which isn't a general-purpose function: React made
it for one record, the record of \texttt{Counter}, the first time
\texttt{Counter} called \texttt{useState}. Calling it marks that record as
needing a new drawing, and React then runs the function the record keeps.
That's also why two counters on the same page never mix their numbers: each has
its own record, its own \texttt{count} and its own \texttt{setCount}.
\end{question}

\section{A component isn't a function you call}

The click hides one more question. \texttt{useState(0)} takes no name, so how
does React know which entry of the record it belongs to? By order. While React runs a component, it keeps
a list of that component's hooks. The first \texttt{useState} call gets the
first entry, the second call gets the second, on every render. This is why hooks
must be called the same number of times, in the same order, every time: never
inside an \texttt{if}, never inside a loop.

Since a component is a function, it's tempting to call it like one, writing
\texttt{\{Counter()\}} instead of \texttt{<Counter />}. Both put the same
button on the screen, until the component has state. Calling a component as a
function breaks the order without any \texttt{if} in sight:

\begin{code}
function Counter() {
  const [count] = useState(0);
  return <span>clicked {count} times</span>;
}

function Panel({ show }) {
  const [title] = useState('panel');
  return <div>{title} {show && Counter()}</div>;
}
\end{code}

\texttt{Counter()} is an ordinary call made while React is running
\texttt{Panel}, so its \texttt{useState} lands in \texttt{Panel}'s list.
Rendered first with \texttt{show} false and then with \texttt{show} true,
\texttt{Panel} suddenly has two hooks instead of one:

\begin{terminal}
React has detected a change in the order of Hooks called by Panel.

   Previous render            Next render
   ------------------------------------------------------
1. useState                   useState
2. undefined                  useState
   ^^^^^^^^^^^^^^^^^^^^^^^^^^^^^^^^^^^^^^^^^^^^^^^^^^^^^^

Error: Rendered more hooks than during the previous render.
\end{terminal}

With \texttt{<Counter />} in the same place, it renders ``panel clicked 0
times''. JSX turns \texttt{<Counter />} into a call to the runtime with
\texttt{Counter} as an argument, not a call to \texttt{Counter}. React decides
when to run it and gives it its own list.

Hooks rely on a rule that nothing in the language enforces. It's an odd design,
and it was the right trade: the code got much shorter, and a lint rule,
\texttt{rules-of-hooks}, checks the order for you. Turn it on, and never call a
component yourself.

\section{The browser decides who you can talk to}

A page loaded from \texttt{localhost:5173} that calls an API on
\texttt{localhost:3000} fails. The API received the request and answered it,
and a \texttt{curl} to the same address works. The page still gets nothing.

\begin{history}{1995--2014 \textperiodcentered{} Same origin}
Netscape's browser had cookies since 1994 and got JavaScript in its version 2,
and together they were dangerous. A script on any page could make the browser
request a page from another site, and the browser would send that site's
cookies along, logged in as you. If the script could read the answer, any page
you visited could read your mail. Netscape's answer, in the same version of the
browser, was the
\textbf{same-origin policy}: a script may read only responses from the same
origin as its page.

Developers spent the next ten years working around it. The best-known trick,
JSONP, which Bob Ippolito described in 2005, loaded data through a
\texttt{<script>} tag, which the policy didn't cover, and only worked because
the page trusted the other site completely. The proper solution started as a
W3C draft in 2006 and became a Recommendation on 16 January 2014:
\textbf{Cross-Origin Resource Sharing}, or CORS, in which the server says, in a
response header, which other origins may read its answers. It now lives in the
WHATWG Fetch standard.
\end{history}

\begin{concept}{Origin}
The scheme, host and port of an address, taken together.
\texttt{http://localhost:5173} and \texttt{http://localhost:3000} are different
origins, and a page opened from a file has the origin \texttt{null}. The
browser sends a cross-origin request, receives the answer, and refuses to hand
it to the script unless the server's headers allow it.

\textbf{What it fixes:} a page you visit can't read what other sites would
show you, logged in as you.

\textbf{What it costs:} a front end and its own API on two ports are already
strangers, and the error appears in the browser, while the server's logs show
a request answered normally.
\end{concept}

There are two ways out. The API can send CORS headers naming the front end's
origin, which means changing the API and getting the headers right for every
route. Or the front end's own server can forward the calls, so the browser
only ever talks to one origin. Vite's dev server has that built in: requests to
a prefix like \texttt{/api} are passed to the real API, and the browser never
learns there's a second server.

\begin{figure}[htbp]
\centering
\begin{tikzpicture}
  \node[hf/reader, minimum width=24mm] (br) at (0,0) {browser\\{\footnotesize page from :5173}};
  \node[hf/bus, minimum width=24mm] (vite) at (5.2,0) {dev server\\{\footnotesize :5173}};
  \node[hf/source, minimum width=24mm] (api) at (10.4,0) {API\\{\footnotesize :3000}};
  \draw[hf/ask] (br) -- node[hf/note, above] {\texttt{/api/health}} (vite);
  \draw[hf/ask] (vite) -- node[hf/note, above] {\texttt{/health}} (api);
  \draw[hf/pain, -{Stealth[length=2mm]}, dashed] (br.south) to[bend right=16] node[pos=0.5, below, hf/note, text=hfAlert] {straight to :3000: the answer arrives, the page can't read it} (api.south);
\end{tikzpicture}
\caption{Above, the browser only talks to the origin that served the page, and
the dev server forwards the call. Below, the direct call that the same-origin
policy blocks.}
\end{figure}

For a front end and an API built together, the proxy is the better answer. Same
origin means no CORS headers to maintain and no extra request the browser sends
first to ask permission. But the proxy exists only while \texttt{npm run dev} is
running. The build produces static files, and whatever serves them in
production has to do the same forwarding.

The proxy's target is a setting, and settings in a front end come with a trap.
Anything the build puts in the bundle is public, because every visitor downloads
it. So Vite only copies into the browser's code the environment variables whose
names start with \texttt{VITE\_}, and leaves everything else on the machine that
runs the build. An address that only the dev server needs should not have that
prefix. A secret should never be in a front end at all, prefix or not.

\section{Testing what the user sees}

A component test runs in Node, where there's no browser and no page. The usual
setup brings in \textbf{jsdom}, an implementation of the DOM written in
JavaScript, and renders the component into it.

What to check is the more interesting question. Airbnb's Enzyme, from 2015, let
tests reach inside a component and inspect its state and its children. Those
tests broke whenever someone reorganized a component, even when the screen
looked exactly the same. In 2018, Kent C.\ Dodds released
\textbf{Testing Library}, built on the opposite rule: find elements the way a
user would, by the text on the screen and the role of the element, and check
what the user would see. A test that finds ``clicked 0 times'', clicks the
button and then finds ``clicked 1 times'' survives any rewrite of the component
that keeps the screen the same.

\section{What it costs}

A ten-line counter ships about 220~kB of JavaScript, and every visitor downloads
it before seeing anything. Almost all of it is React. For a page that will
show a plant diagram, alarms and charts updating every second, that's a fair
price. For a page with one form, it's too much, and plain HTML would do.

The tools change fast. The template that creates a React project with Vite
switched, within a year, to a Rust bundler, a Rust linter as its default and a
TypeScript compiler written in Go, and Create React App, the standard of 2016,
is already gone. A front end is built from more pieces, which move faster than
on the server, and part of the work is keeping up with them.

\section*{Further reading}

\begin{itemize}
  \item P.\ Hunt, ``React: Rethinking best practices'', JSConf EU, 2013.
  \item The React team, ``Our First 50,000 Stars'', 2016,
        \url{https://legacy.reactjs.org/blog/2016/09/28/our-first-50000-stars.html}.
  \item T.\ Occhino and J.\ Chen, ``Hacker Way: Rethinking Web App Development
        at Facebook'', F8, 2014.
  \item D.\ Abramov, ``Why Do React Hooks Rely on Call Order?'', 2018,
        \url{https://overreacted.io/why-do-hooks-rely-on-call-order/}.
  \item J.\ Archibald, ``ECMAScript modules in browsers'', 2017,
        \url{https://jakearchibald.com/2017/es-modules-in-browsers/}.
  \item The React team, ``Sunsetting Create React App'', 2025,
        \url{https://react.dev/blog/2025/02/14/sunsetting-create-react-app}.
  \item A.\ van Kesteren (ed.), \emph{Cross-Origin Resource Sharing}, W3C
        Recommendation, 2014.
  \item K.~C.\ Dodds, ``Introducing the react-testing-library'', 2018.
\end{itemize}
